%% ma: L10-B14-0055
%% lop: 10
%% chuong: 5
%% bai: 14
%% dang: 10.14.9
%% loai: TL
%% muc: VD
%% dap_an: "-"
%% nguon: "chuyen-de-cac-so-dac-trung-cua-mau-so-lieu-khong-ghep-nhom-toan-10-chuong-trinh-moi.pdf (GV. Phan Nhật Linh) (P05), Bài 03 Dạng 2 Bài tập tự luận 15"
%% kien_thuc: "Lớp 10 Bài 13 (số trung bình, trung vị, tứ phân vị, mốt), Bài 14 (khoảng biến thiên, khoảng tứ phân vị, phương sai, độ lệch chuẩn, giá trị bất thường)"
%% trang_thai: chinh-thuc
%% ghi_chu: "Lời giải gốc ghi phương sai Thái Bình 2064,214 và s ≈ 45,43 (sai); đúng s^2 ≈ 2046,21, s ≈ 45,24; kết luận không đổi. Tự duyệt theo yêu cầu giáo viên 22/10/2026."
\begin{ex}
Sản lượng lúa các năm từ 2014 đến 2018 của hai tỉnh Thái Bình và Hậu Giang được cho ở bảng sau (đơn vị: nghìn tấn, nguồn: Tổng cục Thống kê):
\begin{center}
\begin{tabular}{|c|c|c|c|c|c|}\hline
Năm & $2014$ & $2015$ & $2016$ & $2017$ & $2018$\\\hline
Thái Bình & $1061{,}9$ & $1061{,}9$ & $1053{,}6$ & $942{,}6$ & $1030{,}4$\\\hline
Hậu Giang & $1204{,}6$ & $1293{,}1$ & $1231{,}0$ & $1261{,}0$ & $1246{,}1$\\\hline
\end{tabular}
\end{center}
a) Tính độ lệch chuẩn (làm tròn đến hàng phần trăm) và khoảng biến thiên của sản lượng lúa từng tỉnh.

b) Tỉnh nào có sản lượng lúa ổn định hơn? Tại sao?
\loigiai{a) Thái Bình: $\overline x=1030{,}08$; $s^2\approx2046{,}21$; $s\approx45{,}24$; $R=1061{,}9-942{,}6=119{,}3$.

Hậu Giang: $\overline x=1247{,}16$; $s^2\approx875{,}13$; $s\approx29{,}58$; $R=1293{,}1-1204{,}6=88{,}5$.

b) Hậu Giang có độ lệch chuẩn và khoảng biến thiên đều nhỏ hơn nên sản lượng lúa ổn định hơn.}
\end{ex}
